% Vietnamese translation for OI Public Library
% Source-File: IOI中国国家候选队论文/2009/李骥扬/线段跳表——跳表的一个拓展.ppt
\documentclass[11pt]{article}
\usepackage{oipl-vn}

\newcommand{\Oh}{\mathcal{O}}

\title{Danh sách nhảy đoạn: một mở rộng của danh sách nhảy\\{\large Ghi chú theo slide}}
\author{Lý Ký Dương\\Trường Trung học số 2 Thạch Gia Trang}
\date{Luận văn Trại đông 2009}

\begin{document}
\maketitle

\origtitle{Segment skip lists}
\sourcepath{IOI-China-Candidate-Team-Papers/2009/Li-Jiyang/segment-skip-list-slides.ppt}

\section{Mục tiêu}

Slide giới thiệu danh sách nhảy đoạn như một cách mở rộng danh sách nhảy để
xử lý thông tin trên đoạn. Ý tưởng là giữ tính ngẫu nhiên và độ cao logarit
của danh sách nhảy, đồng thời gắn thêm dữ liệu tổng hợp tương tự cây đoạn.

\section{Từ danh sách nhảy đến danh sách nhảy đoạn}

Danh sách nhảy dùng nhiều tầng liên kết để tăng tốc tìm kiếm, chèn và xóa.
Mỗi tầng bỏ qua một số nút ở tầng dưới, nên thao tác từ điển có kỳ vọng
\(\Oh(\log n)\). Với danh sách nhảy đoạn, mỗi liên kết còn đại diện cho một
khoảng phần tử và lưu thông tin của khoảng đó.

\section{Thao tác}

\begin{itemize}
  \item Tìm kiếm đi từ tầng cao xuống thấp như danh sách nhảy thường.
  \item Khi chèn hoặc xóa, các liên kết bị ảnh hưởng phải cập nhật lại thông
  tin đoạn.
  \item Truy vấn đoạn được tách thành một số liên kết phủ rời nhau; kết quả
  được ghép từ thông tin trên các liên kết này.
\end{itemize}

\section{So sánh}

So với cây đoạn, cấu trúc này thuận tiện hơn khi tập khóa thay đổi động. So
với cây cân bằng, nó có cài đặt ngẫu nhiên đơn giản và giữ được trực giác về
các tầng bỏ qua. Bài viết đầy đủ phân tích chi tiết xác suất và chi phí cập
nhật; bản slide nên được đọc để nắm cấu trúc dữ liệu và phạm vi ứng dụng.

\end{document}
